\section{Mechanism: dense adaptive routing and sensor-dropout fragility}
\label{sec:mechanism}

The taxonomy shows GWNet's dense learned routing is fragile under sensor
dropout while AGCRN and STAEformer, which also learn adaptive spatial
relations, are not (Table~\ref{tab:extended}). Two questions follow: is the
dense adaptive adjacency (Eq.~\ref{eq:adaptive_adj}) the driver, and can the
failure mode be mitigated without retraining?

\subsection{Ablation: adaptive vs.\ fixed adjacency}

We run frozen GWNet checkpoints under two inference-time adjacency
configurations. \textbf{GWNet[adaptive]} keeps
$\Aadapt = \text{softmax}(\mathbf{E}_1 \mathbf{E}_2^\top)$;
\textbf{GWNet[fixed]} replaces $\Aadapt$ with the normalized METR-LA distance
adjacency $\Aadj$ and zeros out $\mathbf{E}_1, \mathbf{E}_2$ (rest of the
network unchanged).

\begin{table}[t]
\centering
\caption{\textbf{Replacing GWNet's adaptive adjacency with a fixed road graph halves
sensor-dropout sensitivity and improves city transfer by 22\%, at only 5\% in-distribution
cost.} Sensor dropout: 30\% zero-mask, 5 masks, 3 seeds. City MAE: target-adj, PEMS-BAY (mph),
reported here for interpretability; source-only normalized values in Appendix~\ref{sec:appendix_source_adj}.}
\label{tab:ablation}
\vskip 0.05in
\begin{tabular}{lccc}
\toprule
\textbf{Variant} & \textbf{In-dist MAE} & \textbf{$\DeltaSensor$ (30\%, multi-mask)} & \textbf{City MAE (mph)} \\
\midrule
GWNet[adaptive] & $7.85 \pm 0.11$ & $+0.092 \pm 0.030$ & $4.42 \pm 0.15$ \\
GWNet[fixed]    & $8.27 \pm 0.08$ & $+0.050 \pm 0.010$ & $3.44 \pm 0.04$ \\
\midrule
\emph{Change}   & $+5.3\%$        & $-45.7\%$          & $-22\%$         \\
\bottomrule
\end{tabular}
\end{table}

Fixed substitution cuts sensor-dropout sensitivity by $45.7\%$ and
city-transfer MAE by $22\%$, at $+5.3\%$ in-distribution cost---past the
pre-registered $\geq 30\%$ threshold for attributing fragility to the
adaptive adjacency. Per-horizon decomposition (Table~\ref{tab:consolidated})
shows GWNet's cross-city MAE is nearly horizon-flat ($+4\%$ growth
$h{=}1{\to}12$) while DCRNN grows $+46\%$: the transfer gap is static
structural mis-routing, not error accumulation.

\paragraph{Mask-aware adjacency variant.}
Zeroing and renormalizing adjacency rows/columns for dropped sensors at
inference ($\Aadapt[i,:]{=}\Aadapt[:,i]{=}0$, then
$\Aadapt \leftarrow \Aadapt / \max(\Aadapt\mathbf{1}, \epsilon)$) should
partially restore robustness without retraining.

\begin{table}[t]
\centering
\caption{\textbf{Mask-aware adjacency routing recovers 75\% of GWNet's sensor-dropout gap
without any retraining.} The intervention zeros and renormalizes adjacency rows/columns
for dropped sensors (${\sim}25$~LoC); both variants evaluated on kept sensors only (3 seeds,
5 random 30\% masks each).}
\label{tab:mask_aware}
\vskip 0.05in
\small
\begin{tabular}{lcc}
\toprule
\textbf{Variant} & \textbf{$\DeltaSensor$ (30\%, multi-mask)} & \textbf{Recovery} \\
\midrule
GWNet[standard]    & $+0.092 \pm 0.030$ & --- \\
GWNet[mask-aware]  & $+0.023 \pm 0.018$ & $75.2\%$ \\
\bottomrule
\end{tabular}
\end{table}

Mask-aware recovers $75.2\%$, beating fixed substitution at $45.7\%$:
selectively blocking only the affected rows/columns is better than replacing
the topology with a static prior. The intervention is ${\sim}25$ LoC and
deploys directly on existing GWNet checkpoints.

\subsection{Routing density as a shared mechanism hypothesis}
\label{sec:routing_spectrum}

The GWNet analysis above looks model-specific, but it is the visible tail of a
monotonic relation over \emph{all four} adaptive-graph models in our bank
($n{=}4$ adaptive backbones). Measuring adjacency concentration as
$\text{max\_mass}$ (mean per-row maximum probability) on METR-LA test samples
yields AGCRN $0.008$, MegaCRN $0.008$, MTGNN $0.16$, GWNet $0.23$. Pairing
with each model's $\DeltaSensor$ ($+0.011$, ${-}0.005$, ${+}0.051$,
${+}0.092$) gives Pearson $\rho{=}{+}0.979$ (over the four adaptive
backbones; Appendix~\ref{sec:appendix_routing_spectrum}).
Per-architecture adjacency proxies are canonical and averaged over rows,
samples, and three seeds; STAEformer has no single row-stochastic adjacency and
is excluded from the Pearson (full proxy definitions in
Appendix~\ref{sec:appendix_routing_spectrum}).
The relation replicates on PEMS04 with native checkpoints
($\rho{=}{+}0.910$ on the same four adaptive backbones, vs.\ $+0.979$ on
METR-LA; Appendix~\ref{sec:appendix_pems04_routing}); MegaCRN's
METR-LA-specific negative $\DeltaSensor$ becomes positive on PEMS04
($+0.057$), \emph{consistent with} (rather than fully confirming) the
weekday-prior smoothing diagnosis: a no-covariate retrain remains future
work (\S\ref{sec:discussion}). We treat the result as a two-city replicated
shared-mechanism \emph{hypothesis}, not a settled law.
Within GWNet, top-$k$ sparsification at inference traces the same curve:
relative fragility rises from $0.80\%$ at $k{=}10$ to $1.17\%$ at $k{=}\text{full}$
(Appendix~\ref{sec:appendix_density_sweep}). A first-order spectral bound on
zero-mask perturbation gives per-node corruption $\leq \lvert \mathcal{I} \rvert
\cdot \text{max\_mass}(A)$ up to constants, explaining the linear trend
(Appendix~\ref{sec:appendix_theory}).

\subsection{Ruling out two natural mechanism hypotheses}
\label{sec:polish_ablations}

Table~\ref{tab:mechanism_exclusion} reports ablations that progressively narrow
the cause of MegaCRN and STAEformer's temporal fragility; the final row is an
open residual, not a completed ablation.

\begin{table}[t]
\centering
\small
\caption{\textbf{Mechanism exclusion for MegaCRN/STAEformer temporal fragility.}
``Ruled out'' means the gap is unchanged or worsened after removing the hypothesized driver.}
\label{tab:mechanism_exclusion}
\vskip 0.05in
\begin{tabular}{p{2.6cm}p{3.2cm}p{3.0cm}p{2.5cm}}
\toprule
\textbf{Hypothesis} & \textbf{Test} & \textbf{Result} & \textbf{Verdict} \\
\midrule
TOD/DOW covariate channel
  & Randomize / swap embeddings at inference (5 seeds)
  & MegaCRN: $0\%$ gap change; STAEformer: $12.5\%$ reduction
  & Ruled out (MegaCRN); partial (STAEformer) \\[2pt]
Memory-bank regularization
  & \emph{memzero}: zero memory bank at inference
  & Gap deepens ($-0.075$ vs $-0.040$ at $r{=}0.3$)
  & Ruled out \\[2pt]
Adjacency-mediated smoothing (sensor axis)
  & \emph{attnonly}: identity adjacency
  & $\DeltaSensor$ eliminated across all~$r$
  & Confirmed driver \\[2pt]
Encoder dynamics (temporal axis)
  & ---
  & Open
  & Residual \\
\bottomrule
\end{tabular}
\end{table}

MegaCRN's apparent sensor robustness is adjacency-mediated smoothing toward a
weekday-dominated prior---failure masquerading as robustness
(Appendix~\ref{sec:appendix_memory_ablation}).
The residual temporal fragility lives in the memory bank or encoder dynamics
and is not isolated by inference-time interventions; isolating it requires a
no-covariate retraining, left to future work
(Appendix~\ref{sec:appendix_temporal_cov}).

\subsection{Severity, degree stratification, and imputation baselines}

GWNet is the most sensitive backbone at all five dropout ratios
(GWNet/DCRNN ratio grows $2.5\times \to 4.6\times$ from $10\%$ to $40\%$;
Appendix~\ref{sec:appendix_severity}). Its highest-degree quartile suffers
$\DeltaSensor{=}{+}0.121$ versus $+0.066$ for Q1 ($1.83\times$); DCRNN shows
no degree dependence (Appendix~\ref{sec:appendix_degree_stratification}).
Last-observed-value carry-forward (LOCF) recovers ${\geq}97\%$ of the zero-mask
gap across backbones---the ranking is specific to silent-failure zero-masking
(Appendix~\ref{sec:appendix_imputation}).
